\documentclass{article}
%DIF LATEXDIFF DIFFERENCE FILE
%DIF DEL build/arxiv_changes/old/main.tex   Sun Sep 27 15:53:47 2026
%DIF ADD main.tex                           Sun Sep 27 14:03:18 2026

%DIF 3c3
%DIF < \usepackage{neurips_2026}
%DIF -------
\usepackage[preprint]{neurips_2026} %DIF > 
%DIF -------

\usepackage[utf8]{inputenc}
\usepackage[T1]{fontenc}
\usepackage{hyperref}
\usepackage{url}
\usepackage{booktabs}
\usepackage{amsfonts}
\usepackage{amsmath}
\usepackage{amssymb}   % \lesssim in Appendix A; amsmath alone does not define it
\usepackage{nicefrac}
\usepackage{microtype}
\usepackage{xcolor}
\usepackage{graphicx}
\usepackage{float}     % [H] for the Appendix C figure, so it stays with its paragraph
%DIF 18-20d18
%DIF < \setlength{\abovecaptionskip}{4pt}   % float spacing tightened to hold the body on four pages
%DIF < \setlength{\textfloatsep}{10pt}
%DIF < \setlength{\floatsep}{8pt}
%DIF -------

%DIF 22-29d19
%DIF < % The workshop prescribes the first-page footer text. The style file is otherwise untouched.
%DIF < \makeatletter
%DIF < \renewcommand{\@noticestring}{%
%DIF <   Submitted to the 9th Workshop on Machine Learning and the Physical Sciences (ML4PS 2026).
%DIF <   Do not distribute.%
%DIF < }
%DIF < \makeatother
%DIF < 
%DIF -------
\title{What does self-supervision add to engineered features\\for stellar light curves?}

%DIF 32a21-24
% Author block (roadmap 2.1). Order matches the OpenReview record (Zijun Liu, Yue Ma, Christopher %DIF > 
% Theissen); Yue Ma's affiliation is her OpenReview profile's current entry. Co-author e-mail %DIF > 
% addresses are not printed without their consent. "Hal\i c\i o\u{g}lu" is spelled with the escapes %DIF > 
% rather than raw UTF-8 so the file stays ASCII. %DIF > 
%DIF -------
\author{\DIFdelbegin \DIFdel{Anonymous Author(s)}\DIFdelend %DIF > 
  \DIFaddbegin \DIFadd{Zijun Liu}\DIFaddend \\
  \DIFdelbegin \DIFdel{Affiliation}\DIFdelend \DIFaddbegin \DIFadd{University of California, San Diego}\DIFaddend \\
  \DIFdelbegin \DIFdel{Address}\DIFdelend \DIFaddbegin \DIFadd{\texttt{brl055@ucsd.edu}}\DIFaddend \\
  \DIFdelbegin \DIFdel{\texttt{email}}\DIFdelend \DIFaddbegin \DIFadd{\texttt{liuzijun6688@gmail.com}
  }\And
  \DIFadd{Yue Ma}\\
  \DIFadd{Hal\i c\i o\u{g}lu Data Science Institute}\\
  \DIFadd{University of California, San Diego
  }\And
  \DIFadd{Christopher A. Theissen}\\
  \DIFadd{Department of Astronomy and Astrophysics}\\
  \DIFadd{University of California, San Diego
}\DIFaddend }

%DIF PREAMBLE EXTENSION ADDED BY LATEXDIFF
%DIF UNDERLINE PREAMBLE %DIF PREAMBLE
\RequirePackage[normalem]{ulem} %DIF PREAMBLE
\RequirePackage{color}\definecolor{RED}{rgb}{1,0,0}\definecolor{BLUE}{rgb}{0,0,1} %DIF PREAMBLE
\providecommand{\DIFaddtex}[1]{{\protect\color{blue}\uwave{#1}}} %DIF PREAMBLE
\providecommand{\DIFdeltex}[1]{{\protect\color{red}\sout{#1}}} %DIF PREAMBLE
%DIF SAFE PREAMBLE %DIF PREAMBLE
\providecommand{\DIFaddbegin}{} %DIF PREAMBLE
\providecommand{\DIFaddend}{} %DIF PREAMBLE
\providecommand{\DIFdelbegin}{} %DIF PREAMBLE
\providecommand{\DIFdelend}{} %DIF PREAMBLE
\providecommand{\DIFmodbegin}{} %DIF PREAMBLE
\providecommand{\DIFmodend}{} %DIF PREAMBLE
%DIF FLOATSAFE PREAMBLE %DIF PREAMBLE
\providecommand{\DIFaddFL}[1]{\DIFadd{#1}} %DIF PREAMBLE
\providecommand{\DIFdelFL}[1]{\DIFdel{#1}} %DIF PREAMBLE
\providecommand{\DIFaddbeginFL}{} %DIF PREAMBLE
\providecommand{\DIFaddendFL}{} %DIF PREAMBLE
\providecommand{\DIFdelbeginFL}{} %DIF PREAMBLE
\providecommand{\DIFdelendFL}{} %DIF PREAMBLE
%DIF HYPERREF PREAMBLE %DIF PREAMBLE
\providecommand{\DIFadd}[1]{\texorpdfstring{\DIFaddtex{#1}}{#1}} %DIF PREAMBLE
\providecommand{\DIFdel}[1]{\texorpdfstring{\DIFdeltex{#1}}{}} %DIF PREAMBLE
\providecommand{\DIFscaledelfig}{0.5}
%DIF HIGHLIGHTGRAPHICS PREAMBLE %DIF PREAMBLE
\RequirePackage{settobox} %DIF PREAMBLE
\RequirePackage{letltxmacro} %DIF PREAMBLE
\newsavebox{\DIFdelgraphicsbox} %DIF PREAMBLE
\newlength{\DIFdelgraphicswidth} %DIF PREAMBLE
\newlength{\DIFdelgraphicsheight} %DIF PREAMBLE
% store original definition of \includegraphics %DIF PREAMBLE
\LetLtxMacro{\DIFOincludegraphics}{\includegraphics} %DIF PREAMBLE
\providecommand{\DIFaddincludegraphics}[2][]{{\color{blue}\fbox{\DIFOincludegraphics[#1]{#2}}}} %DIF PREAMBLE
\providecommand{\DIFdelincludegraphics}[2][]{% %DIF PREAMBLE
\sbox{\DIFdelgraphicsbox}{\DIFOincludegraphics[#1]{#2}}% %DIF PREAMBLE
\settoboxwidth{\DIFdelgraphicswidth}{\DIFdelgraphicsbox} %DIF PREAMBLE
\settoboxtotalheight{\DIFdelgraphicsheight}{\DIFdelgraphicsbox} %DIF PREAMBLE
\scalebox{\DIFscaledelfig}{% %DIF PREAMBLE
\parbox[b]{\DIFdelgraphicswidth}{\usebox{\DIFdelgraphicsbox}\\[-\baselineskip] \rule{\DIFdelgraphicswidth}{0em}}\llap{\resizebox{\DIFdelgraphicswidth}{\DIFdelgraphicsheight}{% %DIF PREAMBLE
\setlength{\unitlength}{\DIFdelgraphicswidth}% %DIF PREAMBLE
\begin{picture}(1,1)% %DIF PREAMBLE
\thicklines\linethickness{2pt} %DIF PREAMBLE
{\color[rgb]{1,0,0}\put(0,0){\framebox(1,1){}}}% %DIF PREAMBLE
{\color[rgb]{1,0,0}\put(0,0){\line( 1,1){1}}}% %DIF PREAMBLE
{\color[rgb]{1,0,0}\put(0,1){\line(1,-1){1}}}% %DIF PREAMBLE
\end{picture}% %DIF PREAMBLE
}\hspace*{3pt}}} %DIF PREAMBLE
} %DIF PREAMBLE
\LetLtxMacro{\DIFOaddbegin}{\DIFaddbegin} %DIF PREAMBLE
\LetLtxMacro{\DIFOaddend}{\DIFaddend} %DIF PREAMBLE
\LetLtxMacro{\DIFOdelbegin}{\DIFdelbegin} %DIF PREAMBLE
\LetLtxMacro{\DIFOdelend}{\DIFdelend} %DIF PREAMBLE
\DeclareRobustCommand{\DIFaddbegin}{\DIFOaddbegin \let\includegraphics\DIFaddincludegraphics} %DIF PREAMBLE
\DeclareRobustCommand{\DIFaddend}{\DIFOaddend \let\includegraphics\DIFOincludegraphics} %DIF PREAMBLE
\DeclareRobustCommand{\DIFdelbegin}{\DIFOdelbegin \let\includegraphics\DIFdelincludegraphics} %DIF PREAMBLE
\DeclareRobustCommand{\DIFdelend}{\DIFOaddend \let\includegraphics\DIFOincludegraphics} %DIF PREAMBLE
\LetLtxMacro{\DIFOaddbeginFL}{\DIFaddbeginFL} %DIF PREAMBLE
\LetLtxMacro{\DIFOaddendFL}{\DIFaddendFL} %DIF PREAMBLE
\LetLtxMacro{\DIFOdelbeginFL}{\DIFdelbeginFL} %DIF PREAMBLE
\LetLtxMacro{\DIFOdelendFL}{\DIFdelendFL} %DIF PREAMBLE
\DeclareRobustCommand{\DIFaddbeginFL}{\DIFOaddbeginFL \let\includegraphics\DIFaddincludegraphics} %DIF PREAMBLE
\DeclareRobustCommand{\DIFaddendFL}{\DIFOaddendFL \let\includegraphics\DIFOincludegraphics} %DIF PREAMBLE
\DeclareRobustCommand{\DIFdelbeginFL}{\DIFOdelbeginFL \let\includegraphics\DIFdelincludegraphics} %DIF PREAMBLE
\DeclareRobustCommand{\DIFdelendFL}{\DIFOaddendFL \let\includegraphics\DIFOincludegraphics} %DIF PREAMBLE
%DIF AMSMATHULEM PREAMBLE %DIF PREAMBLE
\makeatletter %DIF PREAMBLE
\let\sout@orig\sout %DIF PREAMBLE
\renewcommand{\sout}[1]{\ifmmode\text{\sout@orig{\ensuremath{#1}}}\else\sout@orig{#1}\fi} %DIF PREAMBLE
\makeatother %DIF PREAMBLE
%DIF COLORLISTINGS PREAMBLE %DIF PREAMBLE
\RequirePackage{listings} %DIF PREAMBLE
\RequirePackage{color} %DIF PREAMBLE
\lstdefinelanguage{DIFcode}{ %DIF PREAMBLE
%DIF DIFCODE_UNDERLINE %DIF PREAMBLE
  moredelim=[il][\color{red}\sout]{\%DIF\ <\ }, %DIF PREAMBLE
  moredelim=[il][\color{blue}\uwave]{\%DIF\ >\ } %DIF PREAMBLE
} %DIF PREAMBLE
\lstdefinestyle{DIFverbatimstyle}{ %DIF PREAMBLE
	language=DIFcode, %DIF PREAMBLE
	basicstyle=\ttfamily, %DIF PREAMBLE
	columns=fullflexible, %DIF PREAMBLE
	keepspaces=true %DIF PREAMBLE
} %DIF PREAMBLE
\lstnewenvironment{DIFverbatim}[1][]{\lstset{style=DIFverbatimstyle}}{} %DIF PREAMBLE
\lstnewenvironment{DIFverbatim*}[1][]{\lstset{style=DIFverbatimstyle,showspaces=true}}{} %DIF PREAMBLE
\lstset{extendedchars=\true,inputencoding=utf8}

%DIF END PREAMBLE EXTENSION ADDED BY LATEXDIFF

\begin{document}

\maketitle

\begin{abstract}
TESS light curves encode a star's pulsation, rotation, binarity and oscillations, among other physical
properties, but the survey
produces far more light curves than can be labelled by hand. Astronomers have engineered a few
dozen physically motivated features from them, which together remain a strong baseline. Previous work has trained
encoders on light curves and judged them by their own downstream scores, but has not asked what
they add to the engineered features. We pre-train a convolutional encoder on unlabelled TESS light curves to predict each
window's latent representation from those of its neighbours. We then freeze it, add its features to
the 25 engineered features, and read the combination out with one fixed linear probe on ten tasks:
five variability classes, four asteroseismic targets and rotation period. Adding the learned features
improves the readout on 7 of 10 tasks over the engineered features alone, while an untrained encoder
does not. The combination also outperforms two networks trained end-to-end on the same labelled
stars on most tasks. An encoder trained without labels therefore gives astronomers new features
that complement the ones they already use rather than replace them.
\end{abstract}

\section{Introduction}

A star's light curve carries its pulsation, rotation, eclipses and oscillations, and with them its
age, interior and companions. Wide-field surveys now produce light curves far faster
than they can be labelled. TESS
\citep{ricker2015} returns 2-minute or 20-second \citep{huber2022} light curves for hundreds of
thousands of pre-selected stars, and full-frame images of every star in the field at cadences that
change between mission cycles (30 minutes, then 10, now 200 seconds). LSST \citep{ivezic2019}
will add \emph{billions} more light curves, albeit at a cadence of days. Self-supervised learning is the
natural response: pre-train one encoder on the unlabelled archive, then attach a cheap readout for
each task. Encoders of this shape exist for ground-based, Gaia and Kepler light curves
\citep{donoso2025astromer2,zuo2025falco,rui2026jepa} and are evaluated by classification accuracy,
sometimes against engineered features \citep{rui2026jepa}; a label-free TESS representation
\citep{poznanski2026} is evaluated by whether repeat observations of one star lie close in its embedding space.
None measures what they add to engineered features.

In this domain, the baseline that a learned encoder must improve on is already strong. Decades of variable-star
research have distilled light curves into compact engineered summaries: periodogram peak structure,
band-limited power, spectral entropy and robust scatter statistics. Pipelines built on them remain
highly competitive. The benchmark of \citet{li2025staremb} makes the point
sharply: across seven expert-labelled variable-star classes, no time-series foundation model they
tested strictly surpassed the long-established domain feature baseline on classification. A learned
encoder for this data therefore has to answer one question: do its features add predictive signal to
the engineered features under the same readout?

We measure that directly. We pre-train an encoder on TESS light curves, freeze it, and add its
features to the 25 engineered features under one fixed linear readout. We ask two things: what the
learned features add, and whether the combination, which trains only a linear model on the labels,
outperforms two networks trained end-to-end on the same labels. Our contributions are:

\begin{itemize}
\item \textbf{The learned features add to the engineered features rather than replacing them.}
Adding the learned features to the engineered ones improves a fixed linear readout beyond
$2\,\mathrm{SE}$ on 7 of 10 tasks; an untrained encoder gives no such gain, so the gain comes from pre-training,
not from merely adding 128 columns.
\item \textbf{The combined pipeline, read out by a fixed linear model, outperforms the two
supervised networks we train end-to-end on the same labels on most tasks.} Against a supervised Conv1D with the
encoder's architecture it is ahead on 6
of 10 tasks and indistinguishable on the other 4; against an MLP on raw flux it is ahead on 9 of 10.
\end{itemize}

\section{Data and method}

\paragraph{Data.} We use 2-minute SPOC PDCSAP light curves of stars brighter than TESS magnitude
$T = 10$ from Sectors 27--101. Pre-training and five of the ten probes (pulsating, eclipsing binary,
rotation, transit and rotation period) use a label-enriched subset of 13{,}470 stars, split 70/15/15
by star: every star labelled as a transit host, eclipsing binary or pulsator that has a usable
window, plus a random sample of 10{,}000 quiet stars, listed in none of our variability catalogues. The pre-training
loss never sees a label, but this pool is label-enriched, not a random draw from the survey. The other five probes draw their labels from catalogues that barely overlap this subset and are scored on a second pool of 22{,}860 stars built the same way (every catalogue positive
with a usable window plus 10{,}000 quiet stars, split 70/15/15). No test star of the first pool
enters pre-training. \DIFaddbegin \DIFadd{The first pool admits positives only through the transit, eclipsing-binary and pulsating catalogues and screens its quiet stars against the rotation catalogue as well, so every rotator in it carries a second label and it holds no rotation-only star; Appendix~\ref{app:rotpool} rescores the two rotation tasks on a third pool of 106}{\DIFadd{,}}\DIFadd{284 stars, the corpus outside the first pool, where rotators are a natural 13\%. }\DIFaddend About 8\% of the second pool's test stars (265 of 3{,}429) were in the pre-training subset; removing them leaves all seven gains in place
(Appendix~\ref{app:unseen}). Each light curve is split into contiguous segments wherever the time axis
breaks, at a missing cadence or a jump longer than five median cadences; segments are
median-normalised by their median absolute deviation and sliced into non-overlapping windows of 256
cadences ($256 \times 120$\,s $\approx 8.5$\,h); restricting to one cadence keeps every window the
same span of time. Windows containing a missing value are discarded, never interpolated or padded,
and segments are never stitched across sector or gap boundaries. Pre-training consumes
sequences of exactly 16 consecutive windows ($\approx 5.69$\,d), cut at a random offset within a
segment, and uses no labels.

\paragraph{Encoder.} A 1D convolutional variational encoder \citep{kingma2014} maps each window to a
128-dimensional Gaussian posterior, and a GRU \citep{cho2014} predicts the next window's latent from
its predecessors, run both forward and backward over the sequence. The objective is
\begin{align*}
\mathcal{L} = {} & \|\hat{x}-x\|^{2} + \beta\,\mathrm{KL}
+ \lambda\big[\|\hat{z}_{t+1}-\bar{z}_{t+1}\|^{2} + \|\hat{z}_{t-1}-\bar{z}_{t-1}\|^{2}\big] \\
& + w\big[\|\log P_{H}(\hat{x})-\log P_{H}(x)\|^{2} + \|\Delta\hat{x}-\Delta x\|^{2}\big],
\end{align*}
where every term is a mean squared error, $\bar{z}$ the stop-gradient target of the forward and
backward predictors, $P_{H}$ the Hann-tapered power spectrum, $\Delta$ the first difference, and
$(\beta,\lambda,w) = (0.1, 60, 0.3)$, $\beta$ warmed up linearly over ten epochs. The encoder is trained with six seeds and then frozen; no gradient from any label reaches
it.

\paragraph{Representations and readout.} A star's learned features are $\mu$, the mean, over every
window of its first segment (16 for 95\% of stars, 20 for the rest), of the per-window posterior
means. The engineered baseline is a fixed 25-feature vector modelled on the T'DA variability pipeline
\citep{audenaert2021}, computed on the same first segment, 19 from its periodogram and 6 time-domain statistics; Appendix~\ref{app:features} defines each feature. Readouts are fixed in advance and never
tuned on a task, with inputs standardised using training-split statistics alone: class-balanced $L_{2}$ logistic regression
at $C = 1$ for the detection tasks, ridge for the two regressions with its penalty set by leave-one-out
cross-validation. We compare four
arms, each the same readout on a different feature set: \texttt{features}, $\mu$,
\texttt{features}\,$\oplus\,\mu$, and, as a control, the same fusion built on an \emph{untrained}
encoder; the two supervised baselines below are the fifth and sixth arms.
Appendix~\ref{app:readout} repeats the first three under XGBoost.

\paragraph{Supervised baselines.} Two end-to-end baselines see exactly the same input windows, star
sets and splits. \textbf{C1} reuses the encoder's convolutional trunk unchanged, adds a 128-unit linear layer
and a scalar output, and is trained per task with a class-weighted loss on a star-level score, the mean of the star's window
logits computed inside the network so that the gradient flows through the mean. \textbf{C2} replaces the convolutional trunk with a
dense one and changes nothing else, so C1 versus C2 isolates the convolutional inductive bias. Both
use one pre-registered regularisation setting for all tasks, select on a validation split no probe
has ever touched, and are trained with six seeds.

\paragraph{Tasks and metrics.} We score ten tasks with one probe per physical quantity: pulsating
variables \citep{gao2025}, eclipsing binaries \citep{prsa2022}, rotation and rotation period
\citep{boyle2026}, transits \citep{guerrero2021}, flare activity \citep{seli2025,vida2021},
oscillating red giants and their $\nu_{\max}$ \citep{hon2021}, solar-like oscillators
\citep{hatt2023}, and RGB versus helium-burning discrimination \citep{sreenivas2026}. Here
$\nu_{\max}$ is the frequency of maximum oscillation power. Detection tasks are scored by the area
under the precision--recall curve (PR-AUC), whose floor is the positive-class prevalence rather than
$0.5$; prevalence is therefore quoted with every detection score. The two regressions are scored by
$R^{2}$.

\paragraph{Uncertainties.} All error bars are $2\times$ the standard error of the mean over seeds.
For the $\mu$ arms the seeds are independent pre-training runs (six); for C1 and C2 they are
initialisation and shuffling seeds (six); the engineered arm is
deterministic. The seeds of the $\mu$ arms and of the supervised arms are unrelated, so differences
between them are treated as unpaired and their bands combined as $\sqrt{\mathrm{SE}_a^{2}+\mathrm{SE}_b^{2}}$.
A difference counts only when it exceeds the combined $2\,\mathrm{SE}$ band. These bands
measure re-training spread only; a paired bootstrap over the test stars (1{,}000 resamples, seeds
averaged; Appendix~\ref{app:bootstrap}) checks that no verdict hinges on the particular test stars.

\section{Results}

\begin{table}[t]
\caption{Ten-task scorecard. Frozen arms: a fixed logistic or ridge readout, or XGBoost on the same
inputs (Appendix~\ref{app:readout}). Bold: best point estimate of the three arms under each frozen
readout\DIFaddbeginFL \DIFaddFL{, with no significance test attached; bold in the appendix tables instead marks a difference
resolved beyond its error bar, so the two conventions are not interchangeable}\DIFaddendFL . $n$: test stars;
prev.: positive-class prevalence.}
\label{tab:scorecard}
\centering
\scriptsize
% GENERATED by experiments/plot_ml4ps_scorecard.py -- do not hand-edit
\setlength{\tabcolsep}{3pt}
\begin{tabular}{llrrrrrrrrr}
\toprule
 & & & \multicolumn{3}{c}{frozen linear readout} & \multicolumn{3}{c}{frozen XGBoost readout} & \multicolumn{2}{c}{end-to-end supervised} \\
\cmidrule(lr){4-6}\cmidrule(lr){7-9}\cmidrule(lr){10-11}
task & metric & $n$ (prev.) & features & $\mu$ & features\,$\oplus$\,$\mu$ & features & $\mu$ & features\,$\oplus$\,$\mu$ & Conv1D & MLP \\
\midrule
pulsating & PR-AUC & 2021 (0.107) & 0.789 & 0.806 & \textbf{0.846} & \textbf{0.845} & 0.821 & 0.834 & 0.816 & 0.748 \\
eclipsing binary & PR-AUC & 2021 (0.097) & 0.742 & 0.771 & \textbf{0.775} & 0.798 & 0.789 & \textbf{0.807} & 0.774 & 0.719 \\
rotation & PR-AUC & 2021 (0.089) & 0.540 & \textbf{0.559} & 0.555 & \textbf{0.596} & 0.544 & 0.576 & 0.565 & 0.526 \\
transit & PR-AUC & 2021 (0.060) & 0.190 & 0.144 & \textbf{0.218} & 0.217 & 0.223 & \textbf{0.228} & 0.134 & 0.109 \\
oscillating giant & PR-AUC & 3429 (0.383) & \textbf{0.920} & 0.854 & 0.918 & \textbf{0.924} & 0.899 & 0.923 & 0.875 & 0.883 \\
solar-like osc. & PR-AUC & 3429 (0.122) & 0.320 & 0.336 & \textbf{0.390} & \textbf{0.470} & 0.409 & 0.463 & 0.396 & 0.275 \\
RGB vs. HeB & ROC-AUC & 161 (0.702) & \textbf{0.758} & 0.660 & 0.708 & 0.746 & 0.725 & \textbf{0.763} & 0.530 & 0.762 \\
$\nu_{\max}$ & $R^2$ & 1313 & 0.831 & 0.802 & \textbf{0.867} & 0.916 & 0.831 & \textbf{0.919} & 0.832 & 0.834 \\
rotation period & $R^2$ & 150 & 0.703 & 0.677 & \textbf{0.717} & 0.695 & 0.655 & \textbf{0.730} & 0.597 & 0.698 \\
flare activity & PR-AUC & 3429 (0.089) & 0.474 & 0.453 & \textbf{0.524} & 0.562 & 0.485 & \textbf{0.565} & 0.533 & 0.444 \\
\bottomrule
\end{tabular}

\end{table}

\begin{figure}[t]
\centering
\includegraphics[width=\linewidth]{figures/fig1_deltas_xgb.pdf}
\caption{The two claims. \textbf{(a)} What the learned features add to the engineered features, under
the fixed linear readout of Table~\ref{tab:scorecard} (solid) and under an XGBoost readout on the
same inputs (faded; Appendix~\ref{app:readout}). \textbf{(b)} The fusion probe against the two
supervised baselines. Bars are $2\,\mathrm{SE}$.}
\label{fig:deltas}
\end{figure}

\paragraph{The learned features add to the engineered features.} Table~\ref{tab:scorecard} and
Figure~\ref{fig:deltas}a give the comparison. Adding $\mu$ to the 25 features improves the
fixed readout beyond $2\,\mathrm{SE}$ on 7 of 10 tasks, led by solar-like oscillators $+0.071$,
pulsating $+0.057$ and flare activity $+0.051$, and reaching $+0.013$ ($R^{2}$) on rotation period.
Rotation is positive but just short of the
bar ($+0.016$ against $0.016$). \DIFaddbegin \DIFadd{On the rotation pool of Appendix~\ref{app:rotpool}, where the two classes differ in rotation and nothing else, the rotation gain is $+0.055 \pm 0.006$ but an untrained encoder also gains $+0.029$, so about half of it is what any 128 extra columns buy on 74}{\DIFadd{,}}\DIFadd{000 training stars; the rotation-period gain there is $+0.034 \pm 0.002$ with the untrained control at zero. Neither verdict moves the count. }\DIFaddend The two negative tasks are the ones where 128 extra columns cost more
than they add: the 25 features already cover the task (oscillating giants, $-0.003$ on a baseline
of $0.920$) or the training set is too small to fit 128 more weights (RGB versus helium-burning, $-0.050$ on
755 training stars). On both, an untrained encoder loses by about the same amount as the trained
fusion, within the error band; that is what redundant columns do to a linear fit, not a sign that
$\mu$ carries misleading content (Appendix~\ref{app:negatives}). Across all ten tasks the untrained control is
negative on seven and never exceeds $+0.016$ (Table~\ref{tab:readout}), so the gains are not what any 128 extra columns would
buy on this feature set; they are specific to what pre-training learned. $\mu$ on its own does not
replace the features (ahead on only 4 of 10 tasks alone), as \citet{li2025staremb} also find for
time-series foundation models; we therefore report fusion rather than substitution \citep{makinen2024}.

\paragraph{The fusion probe outperforms the two supervised networks on most tasks.} Against C1, which shares
the convolutional trunk, labels and input windows but is trained end-to-end, the fusion
probe is ahead on 6 of 10 tasks and indistinguishable from it on the other 4
(Figure~\ref{fig:deltas}b). Against
C2, an MLP on raw flux, it is ahead on 9 of 10 and behind on one (RGB versus helium-burning, 161
stars). C1 is not weak: it is ahead of the engineered arm on 5 of 10.
A paired star-bootstrap (Appendix~\ref{app:bootstrap}) excludes zero on 5 of the 7 fusion gains, 5
of the 6 leads over C1 and 7 of the 9 over C2.
This holds for this architecture on 669 to 16{,}002 labelled stars per task, not for
supervision in general.

\paragraph{The GRU prediction term raises the gain on every task.} Removing the
GRU prediction term, all else fixed and the same six
seeds, reduces the fusion gain on all 10 tasks, and by more than $2\,\mathrm{SE}$ on 6 when the two
arms are differenced seed by seed (Table~\ref{tab:readout}). Without it the gain falls by more
than half on 5 of the 8 tasks where fusion gains.

\section{Discussion}

\paragraph{Pick the encoder by a held-out probe, not by its loss.} Across the ten pre-training recipes of
one sweep, a lower validation reconstruction loss went with a \emph{lower} probe score: the Spearman
$\rho$ between negated loss and PR-AUC runs from $-0.58$ to $-0.84$ over four tasks
\DIFaddbegin \DIFadd{(Figure~\ref{fig:valloss_a})}\DIFaddend . The recipe with the
lowest loss is the one without the prediction term on all four, and choosing it would cost
$0.021$--$0.049$ PR-AUC\DIFaddbegin \DIFadd{. On the family our encoder belongs to the anti-correlation does not survive as
such --- it holds on two of the four tasks and reverses on two --- but the loss optimum is still never
the probe optimum }\DIFaddend (Appendix~\ref{app:valloss}). \DIFaddbegin \DIFadd{The same objective carries a second cost: it develops
a localised artifact inside each window, and none of the interventions we tried on four axes cleared
both pre-registered acceptance criteria, so the recipe we ship does so by the stated fallback rule
rather than by winning (Appendix~\ref{app:artifact}). }\DIFaddend An encoder here therefore has to be
chosen by a held-out probe, not by its own loss.

\DIFaddbegin \begin{figure}[t]
\centering
\DIFaddFL{\includegraphics[width=0.85\linewidth]{figures/figB_valloss_a.pdf}
}\caption{\DIFaddFL{Validation reconstruction loss against downstream probe score, one point per pre-training
recipe and one panel per task. The ten recipes share corpus, architecture and loss terms and differ
only in the GRU prediction term (off, forward, forward+backward, multi-step, the last three at three
weights), four seeds each. $x$: best post-warm-up validation reconstruction loss, averaged over the
recipe's seeds; $y$: PR-AUC of the fixed linear probe on $\mu$; $\rho$: Spearman's rank correlation
between negated loss and score. The loss-optimal recipe is the dynamics-free one (red) on all four
tasks, and it is the worst or next-to-worst probe on each. Appendix~\ref{app:valloss} repeats the
measurement on the family this paper's encoder belongs to.}}
\label{fig:valloss_a}
\end{figure}

\DIFaddend \paragraph{\textbf{Limitations and next steps.}} \emph{Readout.} The four linear-readout arms (\texttt{features},
$\mu$, \texttt{features}\,$\oplus\,\mu$, and the untrained control) share one linear model. It has few
parameters, so it cannot rescue weak features, and it is identical across the four, so the comparison
is between feature sets. Our claim is confined to these 25 features and a linear readout: under
XGBoost on the same inputs the gain exceeds $2\,\mathrm{SE}$ on 4 of 10 tasks and is negative on one\DIFaddbegin \DIFadd{;
where the nonlinear readout raises the features-only score, the $\mu$ gain falls, with a rank
correlation of $0.92$ across the ten tasks }\DIFaddend (Appendix~\ref{app:readout}). \emph{Data.} Every arm,
the supervised baselines included, sees one 5.7-day segment of a bright 2-minute star. The fainter,
far more numerous stars in the full-frame images, and light curves spanning several sectors, are the
obvious next experiments. \emph{Labels.} Our
negatives are stars absent from the label catalogues, not stars verified to lack the signal, so
catalogue incompleteness adds positive-unlabelled label noise and the detection scores may be
conservative. The flare label marks stars that have flared somewhere in their TESS data, not a flare inside
the 5.7-day segment we encode, so the task is recognising a flaring star, not finding a flare. \emph{Label scarcity.} When
the readout is trained on a few percent of the labels, the advantage narrows on 7 of 10 tasks, so how
much pre-training helps when labels are scarce remains open. \emph{Discovery.} Every task here has a
catalogue; whether $\mu$ also separates kinds of variability that no catalogue labels yet, for instance
as clusters in the latent space, is untested. \DIFaddbegin \DIFadd{We leave it as an exploratory direction rather than a
result: with no labels there is nothing to score such a separation against, so it needs a design of
its own.
}\DIFaddend 

\paragraph{Reproducibility and disclosure.} \S2 and Appendix~\ref{app:features} give the data,
objective, features and readouts; the encoder has four convolutional stages (widths $32/64/128/256$,
kernel 5, each halving length), a 128-d latent and a one-layer GRU of width 256. Code\DIFdelbegin \DIFdel{and score
tables will be released on publication}\DIFdelend \DIFaddbegin \DIFadd{, configuration
files and the score tables behind Table~\ref{tab:scorecard} and the appendix tables are available at
}\url{https://github.com/BrightonLiu-zZ/stellar-ssl/tree/arxiv-v1}\DIFaddend ; generative AI assisted with
analysis code and prose editing.

\raggedbottom
\bibliographystyle{plainnat}
\begin{thebibliography}{22}
\providecommand{\natexlab}[1]{#1}
\providecommand{\url}[1]{\texttt{#1}}
\expandafter\ifx\csname urlstyle\endcsname\relax
  \providecommand{\doi}[1]{doi: #1}\else
  \providecommand{\doi}{doi: \begingroup \urlstyle{rm}\Url}\fi

\bibitem[Audenaert et~al.(2021)Audenaert, Kuszlewicz, Handberg,
  et~al.]{audenaert2021}
Jeroen Audenaert, James~S. Kuszlewicz, Rasmus Handberg, et~al.
\newblock {TESS} data for asteroseismology ({T'DA}) stellar variability
  classification pipeline: Set-up and application to the {Kepler} {Q9} data.
\newblock \emph{The Astronomical Journal}, 162\penalty0 (5):\penalty0 209,
  2021.
\newblock arXiv:2107.06301.

\bibitem[Bedding et~al.(2011)Bedding, Mosser, Huber, et~al.]{bedding2011}
Timothy~R. Bedding, Benoit Mosser, Daniel Huber, et~al.
\newblock Gravity modes as a way to distinguish between hydrogen- and
  helium-burning red giant stars.
\newblock \emph{Nature}, 471:\penalty0 608, 2011.
\newblock arXiv:1103.5805.

\bibitem[Boyle et~al.(2026)Boyle, Bouma, and Mann]{boyle2026}
Andrew~W. Boyle, Luke~G. Bouma, and Andrew~W. Mann.
\newblock The {TESS} all-sky rotation survey: Periods for 1,046,317 stars
  within 500 pc.
\newblock \emph{arXiv e-prints}, 2026.
\newblock arXiv:2603.05586.

\bibitem[Cho et~al.(2014)Cho, van Merri\"{e}nboer, Gulcehre, et~al.]{cho2014}
Kyunghyun Cho, Bart van Merri\"{e}nboer, Caglar Gulcehre, et~al.
\newblock Learning phrase representations using {RNN} encoder--decoder for
  statistical machine translation.
\newblock In \emph{Conference on Empirical Methods in Natural Language
  Processing}, 2014.
\newblock arXiv:1406.1078.

\bibitem[Donoso-Oliva et~al.(2025)Donoso-Oliva, Becker, Protopapas,
  et~al.]{donoso2025astromer2}
Cristobal Donoso-Oliva, Ignacio Becker, Pavlos Protopapas, et~al.
\newblock {Astromer} 2.
\newblock \emph{arXiv e-prints}, 2025.
\newblock arXiv:2502.02717.

\bibitem[Gao et~al.(2025)Gao, Chen, Wang, and Liu]{gao2025}
Xinyi Gao, Xiaodian Chen, Shu Wang, and Jifeng Liu.
\newblock Classification of periodic variable stars from {TESS}.
\newblock \emph{The Astrophysical Journal Supplement Series}, 276\penalty0
  (2):\penalty0 57, 2025.
\newblock arXiv:2412.06175.

\bibitem[Guerrero et~al.(2021)Guerrero, Seager, Huang, et~al.]{guerrero2021}
Natalia~M. Guerrero, Sara Seager, Chelsea~X. Huang, et~al.
\newblock The {TESS} objects of interest catalog from the {TESS} prime mission.
\newblock \emph{The Astrophysical Journal Supplement Series}, 254\penalty0
  (2):\penalty0 39, 2021.
\newblock arXiv:2103.12538.

\bibitem[Hatt et~al.(2023)Hatt, Nielsen, Chaplin, et~al.]{hatt2023}
Emily Hatt, Martin~B. Nielsen, William~J. Chaplin, et~al.
\newblock A catalogue of solar-like oscillators observed by {TESS} in
  120-second and 20-second cadence.
\newblock \emph{Astronomy \& Astrophysics}, 669:\penalty0 A67, 2023.
\newblock arXiv:2210.09109.

\bibitem[Hon et~al.(2021)Hon, Huber, Kuszlewicz, et~al.]{hon2021}
Marc Hon, Daniel Huber, James~S. Kuszlewicz, et~al.
\newblock A `quick look' at all-sky galactic archeology with {TESS}: 158,000
  oscillating red giants from the {MIT} quick-look pipeline.
\newblock \emph{The Astrophysical Journal}, 919\penalty0 (2):\penalty0 131,
  2021.
\newblock arXiv:2108.01241.

\bibitem[Huber et~al.(2022)Huber, White, Metcalfe, et~al.]{huber2022}
Daniel Huber, Timothy~R. White, Travis~S. Metcalfe, et~al.
\newblock A 20-second cadence view of solar-type stars and their planets with
  {TESS}: Asteroseismology of solar analogs and a re-characterization of pi
  {Men} c.
\newblock \emph{The Astronomical Journal}, 163\penalty0 (2):\penalty0 79, 2022.
\newblock arXiv:2108.09109.

\bibitem[Ivezi\'{c} et~al.(2019)Ivezi\'{c}, Kahn, Tyson, et~al.]{ivezic2019}
\v{Z}eljko Ivezi\'{c}, Steven~M. Kahn, J.~Anthony Tyson, et~al.
\newblock {LSST}: From science drivers to reference design and anticipated data
  products.
\newblock \emph{The Astrophysical Journal}, 873\penalty0 (2):\penalty0 111,
  2019.
\newblock arXiv:0805.2366.

\bibitem[Kingma and Welling(2014)]{kingma2014}
Diederik~P. Kingma and Max Welling.
\newblock Auto-encoding variational bayes.
\newblock In \emph{International Conference on Learning Representations}, 2014.
\newblock arXiv:1312.6114.

\bibitem[Li et~al.(2025)Li, Chen, Rehemtulla, et~al.]{li2025staremb}
Weijian Li, Hong-Yu Chen, Nabeel Rehemtulla, et~al.
\newblock {StarEmbed}: Benchmarking time series foundation models on
  astronomical observations of variable stars.
\newblock \emph{arXiv e-prints}, 2025.
\newblock arXiv:2510.06200.

\bibitem[Makinen et~al.(2024)Makinen, Sui, Wandelt, et~al.]{makinen2024}
T.~Lucas Makinen, Ce~Sui, Benjamin~D. Wandelt, et~al.
\newblock Hybrid summary statistics.
\newblock \emph{arXiv e-prints}, 2024.
\newblock arXiv:2410.07548.

\bibitem[Poznanski(2026)]{poznanski2026}
Dovi Poznanski.
\newblock A useful representation of {TESS} light curves.
\newblock \emph{arXiv e-prints}, 2026.
\newblock arXiv:2605.05324.

\bibitem[Pr\v{s}a et~al.(2022)Pr\v{s}a, Kochoska, Conroy, et~al.]{prsa2022}
Andrej Pr\v{s}a, Angela Kochoska, Kyle~E. Conroy, et~al.
\newblock {TESS} eclipsing binary stars. {I}. short-cadence observations of
  4584 eclipsing binaries in sectors 1--26.
\newblock \emph{The Astrophysical Journal Supplement Series}, 258\penalty0
  (1):\penalty0 16, 2022.
\newblock arXiv:2110.13382.

\bibitem[Ricker et~al.(2015)Ricker, Winn, Vanderspek, et~al.]{ricker2015}
George~R. Ricker, Joshua~N. Winn, Roland Vanderspek, et~al.
\newblock The transiting exoplanet survey satellite.
\newblock \emph{Journal of Astronomical Telescopes, Instruments, and Systems},
  1\penalty0 (1):\penalty0 014003, 2015.
\newblock arXiv:1406.0151.

\bibitem[Rui(2026)]{rui2026jepa}
Yicheng Rui.
\newblock Domain-informed multi-view self-distillation for astronomical
  light-curve representation learning with {JEPA}.
\newblock \emph{arXiv e-prints}, 2026.
\newblock arXiv:2606.28446.

\bibitem[Seli et~al.(2025)Seli, Vida, Ol\'{a}h, et~al.]{seli2025}
B\'{a}lint Seli, Kriszti\'{a}n Vida, Katalin Ol\'{a}h, et~al.
\newblock Stellar flare morphology with {TESS} across the main sequence.
\newblock \emph{Astronomy \& Astrophysics}, 694:\penalty0 A161, 2025.
\newblock arXiv:2412.12989.

\bibitem[Sreenivas et~al.(2026)]{sreenivas2026}
K.~R. Sreenivas et~al.
\newblock Global asteroseismology of 19,000 red giants in the {TESS} continuous
  viewing zones.
\newblock \emph{arXiv e-prints}, 2026.
\newblock arXiv:2604.00498.

\bibitem[Vida et~al.(2021)Vida, B\'{o}di, Szkl\'{e}n\'{a}r, and Seli]{vida2021}
Kriszti\'{a}n Vida, Attila B\'{o}di, Tam\'{a}s Szkl\'{e}n\'{a}r, and B\'{a}lint
  Seli.
\newblock Finding flares in {Kepler} and {TESS} data with recurrent deep neural
  networks.
\newblock \emph{Astronomy \& Astrophysics}, 652:\penalty0 A107, 2021.
\newblock arXiv:2105.11485.

\bibitem[Zuo et~al.(2025)Zuo, Tao, Huang, et~al.]{zuo2025falco}
Xiaoxiong Zuo, Yihan Tao, Yang Huang, et~al.
\newblock {FALCO}: A foundation model of astronomical light curves for
  time-domain astronomy.
\newblock \emph{arXiv e-prints}, 2025.
\newblock arXiv:2504.20290.

\end{thebibliography}


\appendix

\DIFaddbegin \section{\DIFadd{Pre-training and readout details}}
\label{app:training}

\paragraph{\DIFadd{What one pre-training example is.}} \DIFadd{The packed pre-training corpus is the 70\% training
split of the first pool: 9}{\DIFadd{,}}\DIFadd{428 stars, 35}{\DIFadd{,}}\DIFadd{854 segments and 582}{\DIFadd{,}}\DIFadd{352 windows of 256 cadences. The
encoder is monitored on the pool's 15\% validation split (2}{\DIFadd{,}}\DIFadd{021 stars, 121}{\DIFadd{,}}\DIFadd{328 windows), which no
probe ever sees. One example is one segment. Each epoch it yields a fresh run of 16 consecutive
windows beginning at a uniformly random offset inside that segment, so a long segment shows a
different 5.69-day stretch every epoch; a segment with fewer than 16 windows never entered the pack.
Validation uses offset 0 and every segment, so the monitored quantity does not move with the
sampling. A batch is 256 sequences, giving 140 optimiser steps per epoch.
}

\paragraph{\DIFadd{The first-difference term.}} \DIFadd{In the objective of }\S2\DIFadd{, $\Delta x_{t} = x_{t+1} - x_{t}$ is
the first difference along a window, and $\|\Delta\hat{x}-\Delta x\|^{2}$ is the mean squared error
between the reconstruction's cadence-to-cadence changes and the input's. It is there because the
point-wise term $\|\hat{x}-x\|^{2}$ is already small for a smooth curve drawn through a noisy one:
the fastest components of a light curve carry little amplitude, so a decoder that low-passes its
input pays almost nothing for discarding them. Differencing is a high-pass operation, so the same
error re-enters the objective weighted towards high frequency and the smoothed reconstruction is
penalised. $P_{H}$ plays the same role in the frequency domain; the two share the weight $w$ and are
summed with equal weight inside it.
}

\paragraph{\DIFadd{Optimisation.}} \DIFadd{Adam at learning rate $3\times10^{-4}$, cosine-annealed to
$3\times10^{-6}$ over the 100-epoch cap, gradient-norm clipping at $1.0$, mixed precision. $\beta$
rises linearly from $0$ to its target $0.1$ over the first ten epochs and stays there; the shipped
recipe uses no free-bits floor. Training stops early once ten consecutive epochs improve no tracked
checkpoint, a counter that is armed only after the warm-up.
}

\paragraph{\DIFadd{Which checkpoint is frozen.}} \DIFadd{Only post-warm-up epochs are eligible. During the warm-up
$\beta$ and the KL term are still moving, and the scheduled-$\beta$ total is minimised at $\beta=0$,
that is by an untrained model, so an unrestricted argmin can select the warm-up transient rather than
a fit. Among the eligible epochs we freeze the one minimising
$\mathrm{recon} + w\cdot\mathrm{aux}$ on the validation split. The prediction term is deliberately
absent from that metric: it is the quantity a dynamics ablation changes, so including it would select
each arm of that comparison by a different rule.
}

\paragraph{\DIFadd{Readouts.}} \DIFadd{Logistic regression is scikit-learn's }\texttt{\DIFadd{LogisticRegression}} \DIFadd{with $L_{2}$
penalty, $C=1$, }\texttt{\DIFadd{class\_weight="balanced"}} \DIFadd{and at most 2}{\DIFadd{,}}\DIFadd{000 iterations. Ridge is
}\texttt{\DIFadd{RidgeCV}} \DIFadd{over ten log-spaced penalties from $10^{-2}$ to $10^{3}$, chosen by leave-one-out
cross-validation. Every input column is z-scored with the training split's own mean and standard
deviation, and the same two numbers are applied unchanged to validation and test. For the fusion arm
the 25 features and the 128 latent means are column-concatenated first and the 153-column matrix is
standardised as a whole, so no arm gets a different scaling of the columns it shares with another.
}

\paragraph{\DIFadd{The supervised baselines.}} \DIFadd{C1 and C2 are trained with AdamW at learning rate
$3\times10^{-4}$ and weight decay $10^{-4}$, dropout $0.2$ before the task head only (the trunk is
unregularised, as the encoder's is), gradient clipping at $1.0$, mixed precision, batches of 64
stars ($\approx 1{,}040$ windows, since the loss is per star), at most 60 epochs and patience 10.
The positive class enters the loss with weight $n_{\mathrm{neg}}/n_{\mathrm{pos}}$ computed on the
training split. One regularisation setting is fixed in advance for all ten tasks and no per-task
tuning is done. Each run is selected on the best validation score in that task's own headline metric
rather than on validation loss: at the prevalences here the validation-loss minimum and the
validation-metric maximum routinely fall many epochs apart, and selecting on the loss would depress
the baseline we are trying to beat.
}

\section{\DIFadd{The two star pools, and what a quiet star is}}
\label{app:pools}

\DIFadd{Each task is scored on the pool its catalogue populates, and every arm on a task sees the identical
star list, so the two definitions below fix the population but never a comparison.
}

\paragraph{\DIFadd{The first pool (13}{\DIFadd{,}}\DIFadd{470 stars).}} \DIFadd{Every star in the corpus with at least one usable window
that carries a transit, eclipsing-binary or pulsating label: 3}{\DIFadd{,}}\DIFadd{470 stars, assigned to a single
split stratum rarest class first (813 transit, 1}{\DIFadd{,}}\DIFadd{226 eclipsing binary, 1}{\DIFadd{,}}\DIFadd{431 pulsating). To them
are added 10}{\DIFadd{,}}\DIFadd{000 }\emph{\DIFadd{quiet}} \DIFadd{stars drawn uniformly at random, without replacement, from the corpus
stars that appear in none of our five v1 catalogues -- transit, eclipsing binary, pulsating, rotation
and flare. The whole pool is split 70/15/15 by star within each stratum, and the encoder is
pre-trained and monitored on this pool's training and validation splits.
}

\DIFadd{Note what the two clauses do differently. A star becomes a positive member only through the three
detection catalogues; rotation and flare appear only as exclusions on the quiet draw. A rotator that
is not also a transit host, eclipsing binary or pulsator is therefore neither admitted as a positive
nor drawn as a quiet star, and is absent from this pool altogether. Appendix~\ref{app:rotpool} measures what this costs the two rotation rows.
}

\paragraph{\DIFadd{The second pool (22}{\DIFadd{,}}\DIFadd{860 stars).}} \DIFadd{The four external catalogues behind the asteroseismic
and rotation-period tasks are joined to the corpus into one table of 11}{\DIFadd{,}}\DIFadd{007 stars, each row a star
that at least one of them covers. Those 11}{\DIFadd{,}}\DIFadd{007, plus the 1}{\DIFadd{,}}\DIFadd{853 flare-labelled corpus stars not
already among them, are the pool's members (12}{\DIFadd{,}}\DIFadd{860). To them are added a fresh draw of 10}{\DIFadd{,}}\DIFadd{000
quiet stars: corpus stars that carry no transit, eclipsing-binary, pulsating or rotation label, have
no flare label, and are absent from the external table. The split is again 70/15/15 by star, with the
stratum taken rarest first over RGB-versus-helium-burning, Kounkel rotation period, flare, solar-like
oscillator, oscillating giant and quiet.
}

\paragraph{\DIFadd{The two screens are not symmetric.}} \DIFadd{The first pool's quiet draw is screened against the
five v1 catalogues only; it never consults the four external ones, so some of its quiet stars are
oscillating red giants or solar-like oscillators. The second pool's quiet draw is screened against
both sets. Nothing printed here compares a star in one pool with a star in the other, so the
asymmetry costs no verdict; what it does mean is that a negative in either pool is a star
}\emph{\DIFadd{absent from the relevant catalogues}}\DIFadd{, not a star examined and found to lack the signal. The
external table is itself a union of four detection lists, so a star it covers with a zero for
oscillating giants is a star Hatt, Sreenivas or Kounkel listed and \mbox{%DIFAUXCMD
\citet{hon2021} }\hskip0pt%DIFAUXCMD
did not, which is
the same kind of negative. This is the label limitation }\S4 \DIFadd{states, given exactly.
}

\DIFaddend \section{The two tasks where fusion does not help}
\label{app:negatives}

\begin{center}
\small
\begin{tabular}{lrrrrr}
\toprule
task & features & fusion $\Delta$ & $2\,\mathrm{SE}$ & untrained $\Delta$ & $n_{\mathrm{test}}$ \\
\midrule
oscillating giants & 0.920 & $-0.003$ & 0.002 & $-0.003$ & 3{,}429 \\
RGB vs.\ He-burning & 0.758 & $-0.050$ & 0.019 & $-0.038$ & 161 \\
\bottomrule
\end{tabular}
\end{center}

\paragraph{Oscillating giants: the 25 features already cover the task.} The labels of
\citet{hon2021} come from a convolutional classifier reading a $128 \times 128$ image of each star's
power spectrum, which fires on a visible oscillation power excess at
$6 \lesssim \nu_{\max} \lesssim 250\,\mu$Hz. Our engineered vector encodes close to that same
quantity: eight log-spaced band powers spanning 0.3--360\,c/d, plus the spectral centroid and
entropy, locate the excess and describe its concentration. The features accordingly reach $0.920$,
and the trained and untrained fusion arms then lose the same $0.003$: 128 collinear columns
diluting a linear fit, with no trained-versus-untrained difference left to attribute to $\mu$.

\paragraph{RGB versus helium-burning: no arm resolves the discriminant.} \citet{sreenivas2026} assign
the evolutionary state with a convolutional classifier reading the background-corrected power spectrum
\emph{folded at} $\Delta\nu$; the physics underneath is the dipole mixed-mode period spacing
$\Delta\Pi_{1}$ \citep{bedding2011}. Neither is accessible at our baseline. The encoder sees one
256-cadence window at a time, so the finest frequency structure available to $\mu$ is
$\Delta f = 32.6\,\mu$Hz, while $\Delta\nu$ for these stars is $\approx 1.5$--$18\,\mu$Hz; the comb
itself is unresolved, let alone a fold of it. The engineered features, computed on the whole 5.69\,d
segment, reach $2.0\,\mu$Hz, still coarser than the mixed-mode spacing
$\Delta\Pi_{1}\nu^{2} \approx 0.2$--$0.6\,\mu$Hz near $\nu_{\max} = 50\,\mu$Hz. No arm on this row
measures the quantity the label encodes; the $0.758$ the features reach rides on correlates of
evolutionary state rather than on its discriminant. The row behaves accordingly: it is the only one
where C2 ($0.762$) is ahead of the fusion probe while C1 reaches only $0.530$ on 755 training stars,
and it carries the smallest test set in the study.

\DIFaddbegin \paragraph{\DIFadd{Why 128 redundant columns cost a linear fit.}} \DIFadd{A linear readout spends its training stars
estimating one coefficient per column, and each estimate carries its own error into the prediction
whether or not the column carries signal. On the RGB-versus-helium-burning row the features arm fits
25 coefficients on 755 training stars and the fusion arm 153, so the stars-per-coefficient ratio
falls from about 30 to about 5 while, by the resolution argument above, the added columns supply
nothing the label depends on. The variance a column adds scales with how few stars are available to
pin it down; the signal it adds does not. A column is therefore worth its place only when what it
adds exceeds that cost, and on this row and on oscillating giants it does not. The logistic readout's
$L_{2}$ penalty is fixed at $C = 1$ for every task rather than tuned per task, so it damps this effect
but does not remove it. The prediction that follows is the one the table above shows: an untrained
encoder's 128 columns should cost about the same as a trained one's, since the argument never
mentions what the columns contain, and they do ($-0.038$ against $-0.050$, and $-0.003$ against
$-0.003$).
}

\DIFaddend \section{The gain under a nonlinear readout}
\label{app:readout}

\begin{table}[H]
\caption{What 128 added columns change, relative to the 25 features alone, under each frozen readout
(mean over six encoder seeds; absolute scores are in Table~\ref{tab:scorecard}). $+\mu$: our encoder.
No prediction term: the same recipe trained without the GRU prediction term. \DIFaddbeginFL \DIFaddFL{Retained: the share of
the full $\mu$ gain that survives that removal, printed only where the full gain is a gain, and
negative where the gain changes sign rather than shrinking. }\DIFaddendFL Untrained: a randomly
initialised encoder, which is a single encoder and carries no band. Bold: beyond $2\,\mathrm{SE}$ \DIFaddbeginFL \DIFaddFL{---
unlike Table~\ref{tab:scorecard}, where bold marks the best point estimate and implies no test}\DIFaddendFL . The
last row counts tasks won, tied and lost at that bar.}
\label{tab:readout}
\centering
\scriptsize
% GENERATED by experiments/plot_ml4ps_scorecard.py -- do not hand-edit
\DIFdelbeginFL
{\scriptsize\textcolor{red}{\textbf{ML4PS version:}}}\par\nobreak\smallskip
\begingroup\color{red}
\begin{tabular}{llrrrrr}
\toprule
 & & \multicolumn{3}{c}{linear readout} & \multicolumn{2}{c}{XGBoost readout} \\
\cmidrule(lr){3-5}\cmidrule(lr){6-7}
task & metric & $+\mu$ & $+\mu$, no prediction term & $+$untrained $\mu$ & $+\mu$ & $+$untrained $\mu$ \\
\midrule
pulsating & PR-AUC & $\boldsymbol{+}\mathbf{0.057}$ & $\boldsymbol{+}\mathbf{0.051}$ & ${+}0.016$ & ${-}0.011$ & ${-}0.003$ \\
eclipsing binary & PR-AUC & $\boldsymbol{+}\mathbf{0.033}$ & $\boldsymbol{+}\mathbf{0.012}$ & ${-}0.010$ & $\boldsymbol{+}\mathbf{0.010}$ & ${-}0.003$ \\
rotation & PR-AUC & ${+}0.016$ & $\boldsymbol{-}\mathbf{0.011}$ & ${-}0.026$ & $\boldsymbol{-}\mathbf{0.020}$ & ${-}0.034$ \\
transit & PR-AUC & $\boldsymbol{+}\mathbf{0.027}$ & $\boldsymbol{+}\mathbf{0.018}$ & ${+}0.004$ & ${+}0.011$ & ${-}0.027$ \\
oscillating giant & PR-AUC & $\boldsymbol{-}\mathbf{0.003}$ & $\boldsymbol{-}\mathbf{0.006}$ & ${-}0.003$ & ${-}0.001$ & ${-}0.008$ \\
solar-like osc. & PR-AUC & $\boldsymbol{+}\mathbf{0.071}$ & $\boldsymbol{+}\mathbf{0.024}$ & ${-}0.007$ & ${-}0.007$ & ${-}0.026$ \\
RGB vs. HeB & ROC-AUC & $\boldsymbol{-}\mathbf{0.050}$ & $\boldsymbol{-}\mathbf{0.052}$ & ${-}0.038$ & $\boldsymbol{+}\mathbf{0.016}$ & ${-}0.031$ \\
$\nu_{\max}$ & $R^2$ & $\boldsymbol{+}\mathbf{0.036}$ & $\boldsymbol{+}\mathbf{0.010}$ & ${-}0.001$ & $\boldsymbol{+}\mathbf{0.003}$ & ${-}0.005$ \\
rotation period & $R^2$ & $\boldsymbol{+}\mathbf{0.013}$ & ${-}0.005$ & ${-}0.057$ & $\boldsymbol{+}\mathbf{0.036}$ & ${-}0.004$ \\
flare activity & PR-AUC & $\boldsymbol{+}\mathbf{0.051}$ & $\boldsymbol{+}\mathbf{0.033}$ & ${+}0.006$ & ${+}0.003$ & ${+}0.006$ \\
\midrule
win / tie / loss & & 7 / 1 / 2 & & & 4 / 5 / 1 & \\
\bottomrule
\end{tabular}
\endgroup\par\medskip
{\scriptsize\textcolor{blue}{\textbf{arXiv version:}}}\par\nobreak\smallskip
\DIFdelendFL \DIFaddbeginFL \begin{tabular}{llrrrrrr}
\toprule
 & & \multicolumn{4}{c}{linear readout} & \multicolumn{2}{c}{XGBoost readout} \\
\cmidrule(lr){3-6}\cmidrule(lr){7-8}
task & metric & $+\mu$ & $+\mu$, no prediction term & retained & $+$untrained $\mu$ & $+\mu$ & $+$untrained $\mu$ \\
\midrule
pulsating & PR-AUC & $\boldsymbol{+}\mathbf{0.057}$ & $\boldsymbol{+}\mathbf{0.051}$ & $0.89$ & ${+}0.016$ & ${-}0.011$ & ${-}0.003$ \\
eclipsing binary & PR-AUC & $\boldsymbol{+}\mathbf{0.033}$ & $\boldsymbol{+}\mathbf{0.012}$ & $0.38$ & ${-}0.010$ & $\boldsymbol{+}\mathbf{0.010}$ & ${-}0.003$ \\
rotation & PR-AUC & ${+}0.016$ & $\boldsymbol{-}\mathbf{0.011}$ & $-0.69$ & ${-}0.026$ & $\boldsymbol{-}\mathbf{0.020}$ & ${-}0.034$ \\
transit & PR-AUC & $\boldsymbol{+}\mathbf{0.027}$ & $\boldsymbol{+}\mathbf{0.018}$ & $0.66$ & ${+}0.004$ & ${+}0.011$ & ${-}0.027$ \\
oscillating giant & PR-AUC & $\boldsymbol{-}\mathbf{0.003}$ & $\boldsymbol{-}\mathbf{0.006}$ & -- & ${-}0.003$ & ${-}0.001$ & ${-}0.008$ \\
solar-like osc. & PR-AUC & $\boldsymbol{+}\mathbf{0.071}$ & $\boldsymbol{+}\mathbf{0.024}$ & $0.33$ & ${-}0.007$ & ${-}0.007$ & ${-}0.026$ \\
RGB vs. HeB & ROC-AUC & $\boldsymbol{-}\mathbf{0.050}$ & $\boldsymbol{-}\mathbf{0.052}$ & -- & ${-}0.038$ & $\boldsymbol{+}\mathbf{0.016}$ & ${-}0.031$ \\
$\nu_{\max}$ & $R^2$ & $\boldsymbol{+}\mathbf{0.036}$ & $\boldsymbol{+}\mathbf{0.010}$ & $0.29$ & ${-}0.001$ & $\boldsymbol{+}\mathbf{0.003}$ & ${-}0.005$ \\
rotation period & $R^2$ & $\boldsymbol{+}\mathbf{0.013}$ & ${-}0.005$ & $-0.41$ & ${-}0.057$ & $\boldsymbol{+}\mathbf{0.036}$ & ${-}0.004$ \\
flare activity & PR-AUC & $\boldsymbol{+}\mathbf{0.051}$ & $\boldsymbol{+}\mathbf{0.033}$ & $0.65$ & ${+}0.006$ & ${+}0.003$ & ${+}0.006$ \\
\midrule
win / tie / loss & & 7 / 1 / 2 & & & & 4 / 5 / 1 & \\
\bottomrule
\end{tabular}
\DIFaddendFL 

\end{table}

The \DIFaddbegin \DIFadd{retained column is the per-task form of the body's ablation count. On the eight tasks where
fusion gains, removing the GRU prediction term leaves between $-0.69$ and $0.89$ of that gain, below
half on five of them, and on two (rotation and rotation period) the remainder is negative, so the
ablated encoder's columns are worse than not adding them at all. We give the eight fractions rather
than one summary number because the ten gains are measured in PR-AUC, ROC-AUC and $R^{2}$ and cannot
be added together.
}

\DIFadd{The }\DIFaddend XGBoost readout takes the same inputs as the linear one. For each task and seed it is selected
over a grid of maximum depth $\{3,5\}$ and learning rate $\{0.05,0.1\}$, with up to 1{,}000 rounds
and early stopping after 50 rounds without improvement on a held-out 25\% of the training stars;
the test stars play no part in the selection. Under XGBoost the features-only arm also varies with
the seed, through the validation split, so the gain is differenced seed by seed.

Under XGBoost the $\mu$ gain (features\,$\oplus\,\mu$ minus features alone) is smaller than under the
linear readout on most tasks. XGBoost also scores higher than the linear readout on the 25 features
alone on eight of ten tasks, and the two changes go together. On solar-like oscillators the features-only score rises from
0.320 to 0.470 and the $\mu$ gain falls from $+0.071$ to $-0.007$. On RGB versus helium-burning,
where XGBoost does not raise the features-only score, the $\mu$ gain is larger under XGBoost. Across
the ten tasks, the rise in the features-only score and the fall in the $\mu$ gain have a Spearman rank
correlation of 0.92. This is consistent with part of the information $\mu$ supplies to a linear
readout being already present in the 25 features, in a form only a nonlinear readout can use.

Table~\ref{tab:readout} also gives the same comparison with an untrained encoder in place of the
trained one. Under XGBoost its 128 columns lower the score on nine of ten tasks and the trained $\mu$
scores above it on eight, so the gains that survive come from pre-training, not from adding columns.
XGBoost on the 25 features alone also scores above the linear fusion probe on 7 of 10 tasks
(Table~\ref{tab:scorecard}). The body's claim is about what $\mu$ adds under a fixed linear readout,
not that the linear fusion probe is the strongest classifier available on these features.

\section{Validation reconstruction against downstream score}
\label{app:valloss}
\DIFdelbegin %DIFDELCMD < \begin{figure}[H]
%DIFDELCMD < \centering
%DIFDELCMD < %%%
\DIFdelFL{\includegraphics[width=\linewidth]{figures/figB_valloss.pdf}
}%DIFDELCMD < \caption{%
{%DIFAUXCMD
\DIFdelFL{One point per pre-training recipe (mean over its seeds), one panel per task. (a) The
exp05 sweep: ten recipes sharing corpus, architecture and loss terms, differing only in the
GRU prediction term (off, forward, forward+backward, multi-step, the last three at three weights),
four seeds each. (b) The family our encoder belongs to: 19 recipes on the final encoder that
vary the form, weight and floor of the spectral auxiliary term, six seeds each; the diamond has no
auxiliary term, the cross is the one dynamics-off recipe, the star is the recipe we use. $x$: best
post-warm-up validation reconstruction loss; $y$: PR-AUC of the fixed linear probe on $\mu$;
$\rho$: Spearman's rank correlation between negated loss and score.}}
%DIFAUXCMD
%DIFDELCMD < \label{fig:valloss}
%DIFDELCMD < \end{figure}
%DIFDELCMD < %%%
\DIFdelend 

Figure~\DIFdelbegin \DIFdel{\ref{fig:valloss_b} is the measurementbehind the
model-selection paragraph of }%DIFDELCMD < \S4%%%
\DIFdel{.
In row (a)
}\DIFdelend \DIFaddbegin \DIFadd{\ref{fig:valloss_a} in }\S4 \DIFadd{is half of this measurement; Figure~\ref{fig:valloss_b} is the
other half, and the two differ in which family of recipes is varied.
}

\DIFadd{In Figure~\ref{fig:valloss_a} }\DIFaddend the loss-optimal recipe is the dynamics-free one \DIFdelbegin \DIFdel{(red cross) }\DIFdelend on all four tasks, and
it is the worst or next-to-worst probe on each, so selecting on the loss costs $0.021$--$0.049$
PR-AUC against the probe-optimal recipe. The trend is not carried by that one point: with it removed,
$\rho$ still runs from $-0.43$ to $-0.87$. Those recipes are an earlier sweep than our encoder \DIFdelbegin \DIFdel{(}\DIFdelend \DIFaddbegin \DIFadd{--
}\DIFaddend they carry an earlier spectral auxiliary term \DIFdelbegin \DIFdel{), so row (b) }\DIFdelend \DIFaddbegin \DIFadd{-- so Figure~\ref{fig:valloss_b} }\DIFaddend repeats the estimator
on the family we ship: 19 recipes built on the final encoder that vary the form, weight and floor of
the auxiliary term, six seeds each (two more recipes in which one seed collapsed are excluded; we
never drop a seed).
\DIFaddbegin 

\begin{figure}[H]
\centering
\DIFaddFL{\includegraphics[width=\linewidth]{figures/figB_valloss_b.pdf}
}\caption{\DIFaddFL{The same estimator as Figure~\ref{fig:valloss_a}, on the family our encoder belongs to: 19
recipes on the final encoder that vary the form, weight and floor of the spectral auxiliary term, six
seeds each. The diamond has no auxiliary term, the cross is the one dynamics-off recipe, the star is
the recipe we use. $x$: best post-warm-up validation reconstruction loss; $y$: PR-AUC of the fixed
linear probe on $\mu$; $\rho$: Spearman's rank correlation between negated loss and score.}}
\label{fig:valloss_b}
\end{figure}

\DIFaddend There the anti-correlation does not persist as such: $\rho$ is $-0.51$ on pulsating and $-0.50$ on
transit, $+0.37$ on eclipsing binaries and $+0.11$ on rotation. What persists is that the loss is not
a selector: its optimum is again the recipe with the weakest competing term (the diamond, no
auxiliary term at all), which is never the probe optimum, and trusting it costs $0.011$--$0.034$
PR-AUC. So \DIFdelbegin \DIFdel{row (a) }\DIFdelend \DIFaddbegin \DIFadd{the body's figure }\DIFaddend shows an anti-correlation on every task; \DIFdelbegin \DIFdel{in row (b) it survives }\DIFdelend \DIFaddbegin \DIFadd{this one shows it surviving }\DIFaddend on
two tasks and \DIFdelbegin \DIFdel{reverses }\DIFdelend \DIFaddbegin \DIFadd{reversing }\DIFaddend on two; in neither does the loss choose. Both are statements across recipes
only: within one recipe, across its seeds, no label-free curve metric we tried (reconstruction,
reconstruction plus auxiliary, active latent units) ranks the seeds, every such correlation sitting
within $\pm 0.3$ of zero.

\DIFaddbegin \section{\DIFadd{The within-window artifact, and the two gates that governed it}}
\label{app:artifact}

\paragraph{\DIFadd{Why the objective carries a tapered power spectrum.}} \DIFadd{A point-wise reconstruction error
prices a light curve by amplitude, and in these data amplitude sits mostly in the slow component, so
a decoder can score well on it while discarding the periodic structure the tasks depend on. The two
auxiliary terms put the pressure back where the point-wise term does not reach: the first difference
in time, and the log power spectrum across frequency. The spectrum is taken on a Hann-tapered window
rather than on the raw one, for the reason the rest of this appendix gives.
}

\paragraph{\DIFadd{What the decoder found instead.}} \DIFadd{Every term in the objective is computed inside a single
256-cadence window; nothing spans a boundary. A narrow impulse raises log-power at every frequency at
once, so a decoder can lower the log-spectrum error by planting one spike rather than by reproducing
the star's spectrum. Where the cheapest spike sits is decided by the taper, a fixed function of
position within the window: with a rectangular window it is the window edge, and under a Hann taper
it is the window centre, where the taper weight is largest. Figure~\ref{fig:recon} shows both, in
flux and in the position-resolved error. Laid end to end, the spikes form one train with a period of
exactly one window, which is a property of the display and not of the objective.
}

\DIFadd{Measured as the per-position reconstruction error divided by its own interior median (positions
16--240), maximised over position within each seed and then averaged over six seeds, the rectangular
recipe reaches $31.7\times$ at the window edge and the tapered recipe $11.8\times$ near the centre.
The taper moved the defect and shrank it; it did not remove it. An earlier claim that it had was
based on a measurement taken only at the edge, and is withdrawn.
}

\begin{figure}[H]
\centering
\DIFaddFL{\includegraphics[width=\linewidth]{figures/figD_recon.pdf}
}\caption{\DIFaddFL{The artifact, in two stars and in the average. }\textbf{\DIFaddFL{Top two rows:}} \DIFaddFL{three consecutive
256-cadence windows of one pulsator and one quiet star, input flux in grey against each recipe's
reconstruction; dotted lines mark the window boundaries, and both rows of a column share one
$y$-scale. The two recipes differ only in the taper used for the auxiliary spectral term. Neither
decoder reproduces the cadence-scale scatter, which a 128-dimensional bottleneck cannot carry; what
differs is where each one puts a spike it does not need. }\textbf{\DIFaddFL{Bottom:}} \DIFaddFL{reconstruction error by
position within the window, divided by that seed's median over positions 16--240. Thin lines are the
six pre-training seeds, thick lines their mean. The dashed line is the pre-registered artifact gate.
The seed mean understates the tapered recipe's peak because its position moves between seeds
(125--143), which is why the gate is read as a maximum within each seed.}}
\label{fig:recon}
\end{figure}

\paragraph{\DIFadd{The two pre-registered acceptance criteria.}} \DIFadd{Both were registered on 2026-08-15, before
any of the interventions below were written.
}

\emph{\DIFadd{Artifact gate.}} \DIFadd{The within-window impulse ratio -- the quantity plotted in the lower panel of
Figure~\ref{fig:recon} -- taken as the maximum over }\emph{\DIFadd{all}} \DIFadd{positions rather than at the edge
alone, at most $1.5\times$ at full auxiliary weight, over six seeds, }\emph{\DIFadd{and}} \DIFadd{a visual sign-off on
a fresh random sample of reconstructed light curves. The number never passes on its own: reading an
edge-only number is how the withdrawn claim above was made in the first place.
}

\emph{\DIFadd{No-regression gate.}} \DIFadd{No task among eclipsing binaries, pulsators, rotation and transits worse
than the incumbent recipe by more than $2\,\mathrm{SE}$, at both readouts, paired over six seeds. It
is the guard against a clean but useless encoder: a recipe that abolishes the impulse and loses the
probe is not a fix.
}

\paragraph{\DIFadd{What we tried, and why the tapered recipe is the one we ship.}} \DIFadd{We varied the auxiliary
term along four axes: its }\emph{\DIFadd{mechanism}} \DIFadd{(remove it entirely; average the spectrum over a family of
Slepian tapers so that no single position is cheapest; add a penalty on the excess kurtosis of the
residual; both together; put a floor under the spectral sub-term), its }\emph{\DIFadd{amount}} \DIFadd{(six weights of
the best mechanism, from $0.0125$ to $0.20$), the }\emph{\DIFadd{level}} \DIFadd{of that floor (rebuilt as a two-sided
hinge, which unlike a clamp keeps a restoring gradient, at three measured values), and finally a
}\emph{\DIFadd{floor on the KL term}}\DIFadd{, which was the remaining untried lever. No cell cleared both gates. The
ones that removed the impulse lost the probe; the ones that kept the probe kept a measurable impulse.
The nearest miss reached $1.55 \pm 0.05$, which is $2.4\,\mathrm{SE}$ above the gate; its repaired
variant reached $1.47$ over six seeds and $1.55$ over twelve, and regressed on pulsators at both
readouts.
}

\DIFadd{Because no candidate passed, the rule fixed before the ladder ran leaves the incumbent in place. The
encoder this paper uses is therefore the Hann-tapered recipe by a pre-registered fallback, not
because it won a comparison. Two things make that defensible rather than merely procedural. The
artifact was separately measured not to cost the downstream probe -- ablating it moved $\nu_{\max}$
by $+0.0001 \pm 0.0037$ -- and closing the channel, which was also a test of whether the validation
loss becomes a usable selector once it can no longer be lowered by the trick, did not make it one
(Appendix~\ref{app:valloss}).
}

\DIFaddend \section{Paired bootstrap over test stars}
\label{app:bootstrap}

The $2\,\mathrm{SE}$ bands in the body come from re-training and say nothing about whether the
particular test stars favoured one arm. For every printed contrast we therefore resample each task's
test set with replacement 1{,}000 times, rescore every arm on the identical resampled list from its
saved per-star predictions (no model is retrained), average each arm over its seeds, and take the
difference. Table~\ref{tab:bootstrap} gives the 95\% percentile interval of that difference and the
share of resamples in which the fusion arm is ahead; intervals that exclude zero are bold. The check
tightens no verdict and reverses none: every row the body calls ahead has the fusion arm ahead in at
least 65\% of resamples. It does loosen the smallest tasks -- rotation period (150 test stars) and
RGB versus helium-burning (161) account for four of the ten rows where a resolved $2\,\mathrm{SE}$
verdict becomes an interval that touches zero.

\begin{table}[H]
\caption{Paired star-bootstrap of the four contrasts the body counts. Each entry is the 95\% interval
of (left arm $-$ right arm) over 1{,}000 resamples of the test stars, seeds averaged, and the share of
resamples with the difference above zero ($f_{>0}$; 1.00 means the left arm is ahead in every
resample, 0.00 that it is behind in every resample). Bold: interval excludes zero.}
\label{tab:bootstrap}
\centering
\scriptsize
\setlength{\tabcolsep}{2pt}
% GENERATED by experiments/analyze_paper_bootstrap.py -- do not hand-edit
\begin{tabular}{lrrrrrrrr}
\toprule
task & \multicolumn{2}{c}{fusion $-$ features} & \multicolumn{2}{c}{fusion $-$ C1} & \multicolumn{2}{c}{fusion $-$ C2} & \multicolumn{2}{c}{fusion $-$ dynamics-off fusion} \\
 & 95\% CI & $f_{>0}$ & 95\% CI & $f_{>0}$ & 95\% CI & $f_{>0}$ & 95\% CI & $f_{>0}$ \\
\midrule
pulsating & $\mathbf{[+0.012, +0.105]}$ & 0.99 & $[-0.008, +0.073]$ & 0.94 & $\mathbf{[+0.045, +0.150]}$ & 1.00 & $[-0.023, +0.036]$ & 0.69 \\
eclipsing binary & $\mathbf{[+0.005, +0.062]}$ & 0.99 & $[-0.024, +0.025]$ & 0.52 & $\mathbf{[+0.022, +0.094]}$ & 1.00 & $\mathbf{[+0.004, +0.038]}$ & 0.99 \\
rotation & $[-0.026, +0.057]$ & 0.77 & $[-0.056, +0.040]$ & 0.38 & $[-0.023, +0.079]$ & 0.87 & $[-0.005, +0.052]$ & 0.94 \\
transit & $[-0.025, +0.065]$ & 0.85 & $\mathbf{[+0.041, +0.124]}$ & 1.00 & $\mathbf{[+0.046, +0.174]}$ & 1.00 & $[-0.016, +0.035]$ & 0.78 \\
oscillating giants & $[-0.010, +0.005]$ & 0.21 & $\mathbf{[+0.028, +0.057]}$ & 1.00 & $\mathbf{[+0.021, +0.051]}$ & 1.00 & $[-0.002, +0.008]$ & 0.87 \\
solar-like osc. & $\mathbf{[+0.044, +0.097]}$ & 1.00 & $[-0.037, +0.022]$ & 0.34 & $\mathbf{[+0.079, +0.151]}$ & 1.00 & $\mathbf{[+0.025, +0.067]}$ & 1.00 \\
RGB vs.\ HeB & $\mathbf{[-0.087, -0.012]}$ & 0.00 & $\mathbf{[+0.099, +0.247]}$ & 1.00 & $[-0.126, +0.020]$ & 0.08 & $[-0.037, +0.039]$ & 0.52 \\
$\nu_{\max}$ & $\mathbf{[+0.028, +0.043]}$ & 1.00 & $\mathbf{[+0.024, +0.045]}$ & 1.00 & $\mathbf{[+0.018, +0.047]}$ & 1.00 & $\mathbf{[+0.020, +0.031]}$ & 1.00 \\
rotation period & $[-0.038, +0.067]$ & 0.65 & $\mathbf{[+0.054, +0.186]}$ & 1.00 & $[-0.053, +0.100]$ & 0.69 & $[-0.021, +0.060]$ & 0.83 \\
flare & $\mathbf{[+0.019, +0.078]}$ & 1.00 & $[-0.047, +0.029]$ & 0.32 & $\mathbf{[+0.027, +0.128]}$ & 1.00 & $[-0.000, +0.035]$ & 0.97 \\
\bottomrule
\end{tabular}

\end{table}

\section{The 25 engineered features}
\label{app:features}

All 25 features are computed from a star's first segment: its windows concatenated into one flux
series $x$ of 4{,}096 or 5{,}120 cadences, already median-normalised by its median absolute
deviation. The periodogram is the power spectrum of the mean-subtracted series,
$P(f) = |\mathrm{FFT}(x - \bar{x})|^{2}$ with the zero-frequency bin dropped, and
$p(f) = P(f)/\sum_f P(f)$ its normalisation to unit sum. Frequencies are in cycles per day.
$f_{1}, f_{2}, f_{3}$ are the three strongest local maxima of $P$, strongest first.

\begin{table}[H]
\caption{The engineered feature vector, modelled on the T'DA variability pipeline \citep{audenaert2021}.}
\label{tab:features}
\centering
\small
\begin{tabular}{llp{0.56\linewidth}}
\toprule
\# & feature & definition \\
\midrule
1--3 & $\log_{10} f_{1}, \log_{10} f_{2}, \log_{10} f_{3}$ & frequencies of the three strongest periodogram peaks \\
4--6 & $p_{1}, p_{2}, p_{3}$ & normalised power at each peak, $p_{k} = P(f_{k})/\sum_f P(f)$ \\
7--8 & $p_{2}/p_{1}$, $p_{3}/p_{1}$ & peak-power ratios \\
9 & harmonic ratio & $P(2f_{1})/P(f_{1})$ at the bin nearest $2f_{1}$; $0$ if $2f_{1}$ exceeds the Nyquist frequency \\
10--17 & band powers & $\log_{10} \sum_{f \in B_{j}} p(f)$ for eight log-spaced bands $B_{j}$ spanning 0.3--360\,c/d \\
18 & spectral entropy & $-\sum_f p(f)\ln p(f) / \ln N$, $N$ the number of frequency bins \\
19 & spectral centroid & $\log_{10} \sum_f f\,p(f)$ \\
20--21 & skew, kurtosis & sample skewness and excess kurtosis of $x$ \\
22 & point-to-point scatter ratio & $\mathrm{std}(\Delta x)/\mathrm{std}(x)$, $\Delta$ the first difference \\
23 & depth & 95th minus 5th percentile of $x$ \\
24 & MAD & median absolute deviation of $x$ \\
25 & IQR & inter-quartile range of $x$ \\
\bottomrule
\end{tabular}
\end{table}

\DIFaddbegin \paragraph{\DIFadd{What ``modelled on T'DA'' means.}} \DIFadd{\mbox{%DIFAUXCMD
\citet{audenaert2021} }\hskip0pt%DIFAUXCMD
classify Kepler Q9 light curves
into eight variability classes with an ensemble of four supervised learners and a random-forest
metaclassifier. We take none of that: not the classifiers, not the classes, not the training set.
What we borrow is the shape of the feature vector its tree-based members read, and the borrowing is
partial in ways a reader should be able to check.
}

\emph{\DIFadd{Kept.}} \DIFadd{A short list of periodogram peaks with their relative strengths, a harmonic
measure of the dominant peak, an entropy measure of the signal, and robust time-domain scatter
statistics -- skewness, kurtosis and the median absolute deviation, all three of which appear in
their Table~1. Our point-to-point scatter ratio $\mathrm{std}(\Delta x)/\mathrm{std}(x)$ is, up to a
square root, their variability index $\eta_{e}$, the ratio of the mean squared successive difference
to the variance.
}

\emph{\DIFadd{Changed.}} \DIFadd{They extract up to six frequencies with up to ten harmonics each by iterative
prewhitening -- fit a frequency and its harmonics, subtract them, repeat until a significance
criterion stops the loop. We take the three strongest local maxima of one segment's periodogram, with
no prewhitening and no significance test, and keep one harmonic ratio $P(2f_{1})/P(f_{1})$ instead of
a harmonic series. Their amplitudes and phases become, here, powers normalised by the total and
ratios between them: the flux is median-absolute-deviation normalised per segment, so an absolute
amplitude is not a quantity we have, and we drop phases and phase differences entirely. Their
FliPer, the mean power above each of four frequency cut-offs, becomes our eight log-spaced band
powers spanning 0.3--360\,c/d plus a spectral centroid. Their entropy features are computed on the
time series (differential entropy, multiscale sample entropy); ours is the Shannon entropy of the
normalised power spectrum. Because a packed segment is uniformly sampled and gap-free, its
Lomb--Scargle periodogram is the FFT periodogram, so we compute the cheaper one.
}

\emph{\DIFadd{Dropped.}} \DIFadd{Everything that needs a phase fold at a reliably determined period -- the
self-organising-map location, the folded point-to-point statistics and the folded range -- because
5.69 days covers few cycles for the longer-period classes in our menagerie. Also dropped: the
coherency parameter, zero crossings, the range of the cumulative sum, the Shapiro--Wilk statistic,
the variance ratio, the significant-harmonic count and the flux ratio, and the power-spectrum image
their deep-learning member reads, which is not a feature vector at all.
}

\DIFadd{The result is 25 numbers, 19 from the periodogram and 6 from the time domain, and it is a baseline
rather than a reproduction: it is deliberately the cheap, linear-readable part of a pipeline whose
published accuracy came from an ensemble we do not run.
}

\DIFaddend \section{Evaluation restricted to stars never seen in pre-training}
\label{app:unseen}

The encoder is fitted on the training stars of the first pool and monitored on its validation stars,
so the test stars of
the five tasks scored on that pool are unseen by construction. The second pool shares some stars with
it. Table~\ref{tab:unseen} removes every test star that was in the pre-training training or validation
split and rescores the saved per-star predictions; no model is refitted. The gain exceeds
$2\,\mathrm{SE}$ on the same 7 of 10 tasks; the largest change is flare activity, $+0.051$ to
$+0.041$, and the small loss on oscillating giants ($-0.003$) becomes a tie ($-0.001$).

\begin{table}[H]
\caption{The fusion gain (features\,$\oplus\,\mu$ minus features, linear readout) on all test stars and
on the test stars never seen in pre-training. Bold: beyond $2\,\mathrm{SE}$ over six encoder seeds.}
\label{tab:unseen}
\centering
\scriptsize
% GENERATED by experiments/plot_ml4ps_scorecard.py -- do not hand-edit
\begin{tabular}{llrrrr}
\toprule
task & metric & test stars & unseen test stars & gain, all & gain, unseen \\
\midrule
pulsating & PR-AUC & 2021 & 2021 & $\boldsymbol{+}\mathbf{0.057}$ & $\boldsymbol{+}\mathbf{0.057}$ \\
eclipsing binary & PR-AUC & 2021 & 2021 & $\boldsymbol{+}\mathbf{0.033}$ & $\boldsymbol{+}\mathbf{0.033}$ \\
rotation & PR-AUC & 2021 & 2021 & ${+}0.016$ & ${+}0.016$ \\
transit & PR-AUC & 2021 & 2021 & $\boldsymbol{+}\mathbf{0.027}$ & $\boldsymbol{+}\mathbf{0.027}$ \\
oscillating giant & PR-AUC & 3429 & 3164 & $\boldsymbol{-}\mathbf{0.003}$ & ${-}0.001$ \\
solar-like osc. & PR-AUC & 3429 & 3164 & $\boldsymbol{+}\mathbf{0.071}$ & $\boldsymbol{+}\mathbf{0.069}$ \\
RGB vs. HeB & ROC-AUC & 161 & 149 & $\boldsymbol{-}\mathbf{0.050}$ & $\boldsymbol{-}\mathbf{0.048}$ \\
$\nu_{\max}$ & $R^2$ & 1313 & 1218 & $\boldsymbol{+}\mathbf{0.036}$ & $\boldsymbol{+}\mathbf{0.034}$ \\
rotation period & $R^2$ & 150 & 150 & $\boldsymbol{+}\mathbf{0.013}$ & $\boldsymbol{+}\mathbf{0.013}$ \\
flare activity & PR-AUC & 3429 & 3164 & $\boldsymbol{+}\mathbf{0.051}$ & $\boldsymbol{+}\mathbf{0.041}$ \\
\bottomrule
\end{tabular}
\DIFaddbeginFL 

\end{table}

\section{\DIFadd{The two rotation tasks on a pool that contains rotation-only stars}}
\label{app:rotpool}

\paragraph{\DIFadd{The defect.}} \DIFadd{As Appendix~\ref{app:pools} states, the first pool admits a star as a positive
only through the transit, eclipsing-binary and pulsating catalogues, and its quiet draw excludes every
star the rotation catalogue lists. All 1}{\DIFadd{,}}\DIFadd{138 rotators in that pool are therefore also transit hosts,
eclipsing binaries or pulsators (607, 383 and 148 by split stratum), and none of its 10}{\DIFadd{,}}\DIFadd{000 quiet
stars rotates. A readout can separate the two classes of the }\emph{\DIFadd{rotation}} \DIFadd{task by detecting the
co-occurring class, and the }\emph{\DIFadd{rotation period}} \DIFadd{task is fitted on the same co-labelled rotators
(669 training, 150 test). The ML4PS version of this paper scored both tasks on that pool; the body
above still prints those numbers, and this appendix reports the same two tasks on a pool built to
remove the defect. The design, the population, the arms, the readouts and the rule for reading the
result were fixed in writing before any score on the new pool was computed.
}

\paragraph{\DIFadd{The rotation pool.}} \DIFadd{Every corpus star with at least one usable window that is not in the
first pool: 106}{\DIFadd{,}}\DIFadd{284 stars, of which 13}{\DIFadd{,}}\DIFadd{789 carry a rotation period in the TESS All-Sky Rotation
Survey \mbox{%DIFAUXCMD
\citep{boyle2026} }\hskip0pt%DIFAUXCMD
and 92}{\DIFadd{,}}\DIFadd{495 do not. By construction the pool contains no transit host,
eclipsing binary or pulsator (the first pool took every one), no star that entered pre-training, and
no sampling step: rotators are 13.0\% of it, against 12.5\% of the corpus, so both rotation tasks are
scored at the survey's own prevalence rather than the case-control prevalence of the other pools. It
is split 70/15/15 by star, stratified on rotators with period at most five days (8}{\DIFadd{,}}\DIFadd{786), other
rotators (5}{\DIFadd{,}}\DIFadd{003) and non-rotators, seed 0. Rotation period keeps the five-day cap of the body. A
non-rotator here is a star the survey does not list, not a star examined and found not to rotate;
this is the same label limitation that bounds the body's rotation row. No encoder is run again: the
per-star encoder outputs and the 25 engineered features for these stars had been cached for the
survey-prevalence check of }\S4\DIFadd{, and the readouts are the body's, refitted on the new training split.
}

\DIFadd{Two properties still differ between the classes besides rotation, both astrophysical rather than
sampling artefacts: 10.4\% of the rotators and 0.5\% of the non-rotators carry a flare label, and
0.3\% of the rotators against 9.4\% of the non-rotators are oscillating giants or solar-like
oscillators. The headline rows keep every star, because a survey contains them; two sensitivity rows
refit the readout with flare stars removed from both classes, and with the positives restricted to
stars in no non-rotation catalogue and the negatives to stars in no catalogue at all.
}

\paragraph{\DIFadd{Result.}} \DIFadd{Table~\ref{tab:rotpool} gives both tasks on the first pool and on the rotation
pool, every column of Table~\ref{tab:scorecard} included. Three things change and one does not.
The engineered-feature score on the rotation task falls from 0.540 to 0.386: the first pool's
negatives are stars in no catalogue, the rotation pool's negatives are every non-rotator in the
survey, and the second is the harder set. The fusion gain on rotation rises from $+0.016 \pm 0.016$
to $+0.055 \pm 0.006$, but an untrained encoder now also gains $+0.029$ (init spread $+0.027 \pm
0.003$ over six inits) where on the first pool it lost $0.026$; with 74}{\DIFadd{,}}\DIFadd{398 training stars, 128
random columns are no longer redundant to a linear fit. The trained encoder is ahead of the untrained
one by $+0.029 \pm 0.008$, paired by seed. Because the pre-registered rule for these two rows requires
the untrained control to stay within $2\,\mathrm{SE}$ of zero, rotation does not count, as it did
not before. Rotation period does count on both pools, and more clearly on the new one: $+0.034 \pm
0.002$ on 1}{\DIFadd{,}}\DIFadd{318 test stars against $+0.013 \pm 0.010$ on 150, with the untrained control at
$-0.003$. The body's count of 7 of 10 is therefore unchanged. The two sensitivity rows agree with the
headline on both verdicts. The dynamics-off contrast keeps its sign on both tasks ($+0.019 \pm 0.007$
and $+0.020 \pm 0.003$), and a paired star-bootstrap of every contrast in the table
(1}{\DIFadd{,}}\DIFadd{000 resamples, as in Appendix~\ref{app:bootstrap}) agrees with every seed-band verdict.
}

\DIFadd{Two of the body's other tallies would change if the rotation pool replaced the first pool in
Table~\ref{tab:scorecard}, and are stated here so the reader can recount. Under XGBoost the fusion
gain on rotation period turns from $+0.036$ to $-0.004 \pm 0.003$ (the features alone reach 0.829),
so Appendix~\ref{app:readout}'s tally of 4 wins, 5 ties and 1 loss becomes 3, 5 and 2. Against the
raw-flux MLP the fusion probe is behind on rotation period (0.752 against 0.773, $-0.021 \pm 0.004$),
so ``ahead on 9 of 10'' becomes 8 of 10; it stays ahead of the convolutional baseline on both tasks.
Under the linear readout the untrained control's largest gain across the ten tasks becomes the
$+0.029$ above rather than the $+0.016$ the body quotes; that sentence's conclusion, that the gains
which count are not what 128 arbitrary columns buy, is what the rule of this appendix enforces
explicitly.
}

\begin{table}[H]
\caption{\DIFaddFL{The two rotation tasks on the first pool (the body's numbers) and on the rotation pool.
Linear and XGBoost readouts and the two supervised baselines are those of Table~\ref{tab:scorecard};
gain = features\,$\oplus\,\mu$ minus features, mean over six encoder seeds, bold beyond
$2\,\mathrm{SE}$; untrained gain = the same with the untrained encoder. Sensitivity rows refit the
linear readout on the restricted population. $n$: test stars; prev.: rotator prevalence.}}
\label{tab:rotpool}
\centering
\scriptsize
%DIF >  GENERATED by experiments/plot_rotation_pool_appendix.py -- do not hand-edit
\setlength{\tabcolsep}{2.5pt}
\begin{tabular}{llrrrrrrrrrr}
\toprule
 & & & \multicolumn{5}{c}{linear readout} & \multicolumn{2}{c}{XGBoost} & \multicolumn{2}{c}{supervised} \\
\cmidrule(lr){4-8}\cmidrule(lr){9-10}\cmidrule(lr){11-12}
task & population & $n$ (prev.) & features & $\mu$ & fusion & gain & untr.\ gain & features & fusion & Conv1D & MLP \\
\midrule
rotation (PR-AUC) & first pool & 2021 (0.089) & 0.540 & 0.559 & 0.555 & $+0.016$ & $-0.026$ & 0.596 & 0.576 & 0.565 & 0.526 \\
 & rotation pool & 15944 (0.130) & 0.386 & 0.373 & 0.441 & $\mathbf{+0.055}$ & $+0.029$ & 0.518 & 0.515 & 0.409 & 0.420 \\
 & \quad no flare stars & 15670 (0.119) & 0.356 & 0.352 & 0.413 & $\mathbf{+0.057}$ & $+0.029$ & -- & -- & -- & -- \\
 & \quad strict & 14182 (0.129) & 0.367 & 0.372 & 0.427 & $\mathbf{+0.061}$ & $+0.036$ & -- & -- & -- & -- \\
\midrule
rotation period ($R^2$) & first pool & 150 & 0.703 & 0.677 & 0.717 & $\mathbf{+0.013}$ & $-0.057$ & 0.695 & 0.730 & 0.597 & 0.698 \\
 & rotation pool & 1318 & 0.718 & 0.636 & 0.752 & $\mathbf{+0.034}$ & $-0.003$ & 0.829 & 0.825 & 0.697 & 0.773 \\
 & \quad no flare stars & 1174 & 0.713 & 0.628 & 0.744 & $\mathbf{+0.031}$ & $-0.006$ & -- & -- & -- & -- \\
 & \quad strict & 1145 & 0.718 & 0.636 & 0.748 & $\mathbf{+0.031}$ & $-0.004$ & -- & -- & -- & -- \\
\bottomrule
\end{tabular}
\DIFaddendFL 

\end{table}

\end{document}
